\chapter{Configuration Reference}
\label{app:config}

Every configuration setting is loaded from the \file{.env} file in
the repository root through \code{pydantic-settings}
\citep{pydantic}. The complete list is reproduced below and matches
\file{core/config.py}.

\begin{table}[H]
\centering
\caption{Configuration variables (all read from \file{.env}).}
\label{tab:config}
\small
\begin{tabularx}{\linewidth}{@{}lXl@{}}
\toprule
\textbf{Variable} & \textbf{Purpose} & \textbf{Default} \\
\midrule
\code{GROQ\_API\_KEY} & Groq API key. Required for chat/orchestration; if empty, the UI shows an error banner and disables the chat input. & \emph{(empty)} \\
\code{GROQ\_MODEL} & Groq model identifier used by the tool-call loop. & \code{llama-3.3-70b-versatile} \\
\code{MAX\_TOOL\_ITERATIONS} & Maximum number of outer loop iterations per user turn. & \code{8} \\
\code{MAX\_UPLOAD\_MB} & Upload size cap enforced client-side in the Streamlit sidebar. & \code{200} \\
\code{DATA\_DIR} & Root directory for \file{uploads/}, \file{artifacts/}, and \file{tmp/}. Absolute or relative to the repository. & \code{./data} \\
\bottomrule
\end{tabularx}
\end{table}

\section*{Sample \file{.env.example}}

\begin{lstlisting}[style=alBash, caption={Contents of \file{.env.example}.}, label={lst:env}]
# Groq LLM (required for chat/orchestration)
GROQ_API_KEY=gsk_your_key_here
GROQ_MODEL=llama-3.3-70b-versatile

# Orchestrator behavior
MAX_TOOL_ITERATIONS=8
MAX_UPLOAD_MB=200

# Where uploaded datasets, generated artifacts, and tmp files live.
# Leave commented to use ./data at the repo root.
# DATA_DIR=./data
\end{lstlisting}

\section*{Derived Paths}

\code{Settings} exposes the following properties, each of which
creates the underlying directory on first access:

\begin{itemize}[leftmargin=1.4em]
  \item \code{settings.data\_path} $\to$ \file{./data}
  \item \code{settings.uploads\_path} $\to$ \file{./data/uploads}
  \item \code{settings.artifacts\_path} $\to$ \file{./data/artifacts}
  \item \code{settings.tmp\_path} $\to$ \file{./data/tmp}
\end{itemize}
